% Thesis Abstract
\begin{abstract}
  \large \textbf{Post-Landfall Sustainability of Monsoon Depressions
    over India: \\
    The Role of the Brown Ocean Effect, %\\
    Orographic Modulation by the Eastern Ghats, %\\
  and Antecedent Soil Moisture}

  \vspace{0.5cm}

  \normalsize
  % \noindent Monsoon Depressions (MDs) are among the most
  % significant synoptic-scale weather systems of the Indian summer
  % monsoon, accounting for a disproportionate fraction of central
  % and peninsular India's seasonal rainfall. Typically upon
  % landfall, rapid weakening due to the removal of the energy source
  % from ocean and increased surface friction. Yet a subset of MDs
  % exhibits anomalously prolonged inland persistence - an
  % observation inadequately explained by large-scale dynamics alone.
  % This study investigates the land-surface feedback mechanisms that
  % govern this post-landfall sustainability, with a focus on the
  % Brown Ocean Effect (BOE), the modulating role of the Eastern
  % Ghats, and antecedent soil moisture conditions.
  \noindent Monsoon Depressions (MDs) are important synoptic-scale
  systems for India's seasonal rainfall. While most MDs weaken
  rapidly upon landfall due to the loss of oceanic energy and
  increased friction from land interactions, a subset exhibits
  anomalous inland persistence. This study investigates the
  land-surface feedbacks governing this sustainability, focusing on
  the Brown Ocean Effect (BOE), antecedent soil moisture, and the
  orographic influence of the Eastern Ghats.

  \vspace{0.3cm}

  \noindent ERA5 reanalysis data (0.25\(^\circ\) $\times$
  0.25\(^\circ\), hourly) are used as the primary dataset, and
  detected systems are cross-validated against India Meteorological
  Department (IMD) monsoon seasonal reports. MDs are identified using
  850\,hPa relative vorticity thresholds, and their life cycles are
  characterised across multiple Indian summer monsoon seasons using
  track-following diagnostics. Two complementary analytical
  frameworks constitute the primary methodology. Azimuthal
  Hovm\"{o}ller diagrams, constructed by decomposing the MD domain
  into radial annuli at 25\,km intervals extending to 400\,km from
  the storm centre, reveal the spatio-temporal evolution of heat flux
  (surface LHF and SHF), near-surface soil moisture (ERA5
  \texttt{SWVL1}, 0-7\,cm), total precipitation, and 2-metre
  dew-point temperature relative to the MD centre. Track-following
  time series of these variables, stratified by land and ocean
  phases, are used to quantify thermodynamic transitions across the
  land-sea boundary and along the inland track corridor.

  \vspace{0.3cm}

  % \noindent Applying the Andersen and Shepherd (2013) BOE criterion
  % - sustained inland LHF $\geq 70$\,W\,m$^{-2}$ for at least 18
  % consecutive hours - candidate systems are identified and
  % contrasted with those exhibiting conventional post-landfall
  % decay. The analyses are complemented by surface sensible heat
  % flux diagnostics, which document the partitioning of the surface
  % energy budget between latent and sensible heating. Antecedent
  % soil moisture analysis, employing an exponentially weighted index
  % over the future inland track corridor prior to landfall, reveals
  % the role of pre-storm ground wetness in conditioning the land
  % surface for sustained latent heat release. Orographic analysis
  % using high-resolution SRTM elevation data further characterises
  % how the Eastern Ghats barrier modulates translational speed,
  % precipitation organisation, and inland penetration distance.
  \noindent Candidate BOE systems were identified using the Andersen
  and Shepherd (2013) criterion-sustained inland LHF $\geq
  70$\,W\,m$^{-2}$ for $\geq$18 hours and contrasted against
  conventional decaying systems. Antecedent soil moisture was
  assessed using an exponentially weighted index to determine the
  pre-storm ground conditions along future inland tracks.

  \vspace{0.3cm}

  % \noindent Results demonstrate that MDs traversing soils with
  % anomalously high antecedent moisture maintain inland LHF at
  % levels comparable to their ocean phase, sustaining a
  % precipitation-soil moisture positive feedback loop consistent
  % with the BOE mechanism. Systems encountering the high peaks of
  % Eastern Ghats exhibit systematically reduced translational speeds
  % and asymmetric precipitation patterns on windward slopes. The
  % degree of terrain-induced disruption is found to interact with
  % the surface moisture state, suggesting that the BOE can, in part,
  % offset orographically imposed decay for sufficiently
  % pre-moistened land surfaces. These findings advance the
  % mechanistic understanding of why certain MDs exhibit extended
  % inland longevity, with implications for seasonal flood prediction
  % over the Indian subcontinent.
  \noindent Results indicate that MDs traversing regions with high
  antecedent moisture maintain inland LHF levels comparable to the
  ocean phase, sustaining a precipitation-soil moisture feedback loop
  consistent with the BOE. The Eastern Ghats significantly modulate
  these systems, reducing translational speeds and inducing
  asymmetric precipitation on windward slopes. The analysis show that
  a sufficiently pre-moistened land surface can partially offset
  orographically imposed decay. These findings improve the
  mechanistic understanding of MD longevity, providing critical
  insights for seasonal flood forecasting over the Indian subcontinent.

\end{abstract}
